% આ ભાગનો પૂર્ણ સંપાદનીય પાઠ છે. શૈલી, ફોન્ટ, ચિત્રો અને પ્રકરણસંદર્ભ માટે સંપૂર્ણ સ્રોત ZIP તથા reader/full-reader.tex જુઓ. આ એકલી સ્વતંત્ર કમ્પાઇલ થતી મુખ્ય ફાઇલ નથી.

% BEGIN OLP-0516
\olpart{cnt}{પ્રતિતથ્યાત્મક શરતી વિધાનો}
\phantomsection\label{guunit:OLP-0516}


\begingroup
\makeatletter\def\ol@sectioncs{section}\makeatother

% BEGIN OLP-0517
\olchapter{cnt}{int}{પરિચય}
\olchapterid{cnt}{int}
\phantomsection\label{guunit:OLP-0517}


\begingroup
\makeatletter\def\ol@sectioncs{section}\makeatother

% BEGIN OLP-0518
\olfileid{cnt}{int}{mat}

\olsection{ભૌતિક શરતી વિધાન}
\phantomsection\label{guunit:OLP-0518}


અંગ્રેજીમાં શરતી વાક્યનું સૌથી સરળ સ્વરૂપ ``જો \dots તો \dots''
હોય છે. અહીં \dots{}ની જગ્યાએ પોતે વાક્યો આવે છે; દાખલા તરીકે,
``જો બટલરે તે કર્યું હોય, તો માળી નિર્દોષ છે.'' પ્રારંભિક તર્કશાસ્ત્રના
અભ્યાસક્રમોમાં આપણે $\lif$ સંયોજક વડે શરતી વાક્યોનું પ્રતીકીકરણ
શીખીએ છીએ: \dots{}થી દર્શાવેલા ભાગોને, દાખલા તરીકે, સૂત્રો
$!A$ અને~$!B$ વડે પ્રતીકિત કરીએ છીએ; આખું શરતી વાક્ય
$!A \lif !B$ વડે દર્શાવીએ છીએ.

$\lif$ સંયોજક \emph{સત્યતાફલનલક્ષી} છે: $!A \lif !B$નું સત્યમૂલ્ય
---$\True$ કે $\False$---$!A$ અને~$!B$નાં સત્યમૂલ્યો પરથી નક્કી
થાય છે. $!A$ અસત્ય હોય અથવા $!B$ સત્ય હોય તો અને તો જ
$!A \lif !B$ સત્ય છે; અન્યથા તે અસત્ય છે. સત્યમૂલ્ય-નિયોજન
$\pAssign v$ના સંદર્ભમાં, $\pSat/{v}{!A}$ અથવા $\pSat{v}{!B}$ હોય
તો અને તો જ $\pSat{v}{!A \lif !B}$ એમ વ્યાખ્યાયિત કરીએ છીએ. આ
અર્થવિચાર ધરાવતા $\lif$ સંયોજકને \emph{ભૌતિક શરતી વિધાન} કહીએ છીએ.

આ વ્યાખ્યામાંથી કેટલાક પ્રાથમિક તાર્કિક તથ્યો મળે છે. સૌપ્રથમ,
ભૌતિક શરતી વિધાન માટે નિગમન પ્રમેય લાગુ પડે છે:
\begin{align}
  \text{જો } \Gamma, !A \Entails !B \text{ તો } \Gamma \Entails !A \lif !B
\end{align}
તે સત્યતાફલનલક્ષી છે: $!A \lif !B$ અને $\lnot !A \lor !B$ સમતુલ્ય છે:
\begin{align}
  !A \lif !B & \Entails \lnot !A \lor !B\\
  \lnot !A \lor !B & \Entails !A \lif !B
  \intertext{ભૌતિક શરતી વિધાન તેના પરિણામથી તથા તેના પૂર્વપક્ષના નિષેધથી ફલિત થાય છે:}
  !B & \Entails !A \lif !B\\
  \lnot !A & \Entails !A \lif !B
  \intertext{અસત્ય ભૌતિક શરતી વિધાન તેના પૂર્વપક્ષ અને પરિણામના નિષેધના સંયોજનને સમતુલ્ય છે:
    જો $!A \lif !B$ અસત્ય હોય, તો $!A \land \lnot !B$ સત્ય છે અને ઊલટું પણ ખરું:}
  \lnot(!A \lif !B) & \Entails !A \land \lnot !B\\
  !A \land \lnot !B & \Entails \lnot(!A \lif !B)
  \intertext{ભૌતિક શરતી વિધાન પર મોડસ પોનેન્સનો નિયમ લાગુ પડે છે:}
  !A, !A \lif !B & \Entails !B
  \intertext{ભૌતિક શરતી વિધાન પરિણામોને એકત્ર કરે છે:}
  !A \lif !B,  !A \lif !C & \Entails !A \lif (!B \land !C)
  \intertext{પૂર્વપક્ષને હંમેશાં વધુ મજબૂત કરી શકાય છે, એટલે કે આ શરતી વિધાન \emph{એકદિશવર્ધી} છે:}
  !A \lif !B & \Entails (!A \land !C) \lif !B
  \intertext{ભૌતિક શરતી વિધાન સંક્રામક છે, એટલે કે શૃંખલા-નિયમ માન્ય છે:}
  !A \lif !B, !B \lif !C & \Entails !A \lif !C
  \intertext{ભૌતિક શરતી વિધાન તેના પ્રતિપક્ષી સ્વરૂપ સાથે સમતુલ્ય છે:}
  !A \lif !B & \Entails \lnot !B \lif \lnot !A\\
  \lnot !B \lif \lnot !A & \Entails !A \lif !B
\end{align}

ગાણિતિક તર્કવિચારમાં આ બધાં ઉપયોગી અને નિર્વિવાદ અનુમાનો છે. પરંતુ
તત્ત્વચિંતન અને ભાષાવિજ્ઞાનના સાહિત્યમાં, બિનગાણિતિક સંદર્ભોમાં તેમને
અનુરૂપ અનુમાનો સામે મૂકાયેલા પ્રત્યુદાહરણો ઘણાં જોવા મળે છે. તે સૂચવે
છે કે અંગ્રેજીના ``જો \dots તો \dots'' વાક્યોનું પ્રતીકીકરણ કરતી વખતે
ભૌતિક શરતી વિધાન~$\lif$ યોગ્ય સંયોજક નથી---અથવા ઓછામાં ઓછું હંમેશાં
યોગ્ય નથી.
% END OLP-0518
\endgroup


\begingroup
\makeatletter\def\ol@sectioncs{section}\makeatother

% BEGIN OLP-0519
\olfileid{cnt}{int}{par}

\olsection{ભૌતિક શરતી વિધાનના વિરોધાભાસો}
\phantomsection\label{guunit:OLP-0519}


અંગ્રેજીના ``જો \dots તો \dots'' વાક્યોનું પ્રતીકીકરણ કરવા
$!A \lif !B$ વાપરવાની ટીકા સૌપ્રથમ કરનારાઓમાં સી.~આઈ. લૂઈસ હતા.
તેઓ વ્હાઇટહેડ અને રસેલના \emph{પ્રિન્સિપિયા મેથેમેટિકા}માં ભૌતિક
શરતી વિધાનના ઉપયોગની ટીકા કરતા હતા; ત્યાં $\lif$ને ``ફલિત કરે છે''
એમ વાંચવામાં આવતું હતું. લૂઈસે યોગ્ય રીતે વાંધો ઉઠાવ્યો કે જો
$\lif$નો અર્થ ``ફલિત કરે છે'' હોય, તો કોઈ પણ અસત્ય વિધાન~$p$,
$p$માંથી~$q$ ફલિત થાય છે એવું ફલિત કરે છે, કેમ કે $p$ અસત્ય હોય
ત્યારે $p \lif (p \lif q)$ સત્ય છે. એ જ રીતે કોઈ પણ સત્ય વિધાન~$q$,
$p$માંથી~$q$ ફલિત થાય છે એવું ફલિત કરે છે, કેમ કે $q$ સત્ય હોય
ત્યારે $q \lif (p \lif q)$ સત્ય છે.

તર્કશાસ્ત્રીઓ જાણે છે કે \emph{ફલિતતા}, એટલે કે તાર્કિક ફલિતતા,
સંયોજક નથી; તે સૂત્રો અથવા વિધાનો વચ્ચેનો સંબંધ છે. તેથી
ગૂંચવણ ટાળવા $\lif$ને ``ફલિત કરે છે'' એમ વાંચવું ન જોઈએ.\footnote{$\lif$ને
  ``ફલિત કરે છે'' એમ વાંચવાનો પ્રયોગ ગણિતશાસ્ત્રીઓ અને કમ્પ્યુટર
  વૈજ્ઞાનિકોમાં હજુ વ્યાપક છે; જોકે લૂઈસે દર્શાવેલી ગૂંચવણો ટાળવા
  તત્ત્વચિંતકો તેને ``માત્ર ત્યારે જ'' એમ વાંચે છે.} આ ભેદ રાખીએ
તો લૂઈસની ખાસ ચિંતા ઊભી થતી નથી: $p$ને અસત્ય અંગ્રેજી વાક્યનું
પ્રતીક માનીએ તોપણ $p$, $q$ને ``ફલિત'' કરતું નથી. $p \Entails q$
છે કે નહિ તે નક્કી કરવા \emph{બધાં} સત્યમૂલ્ય-નિયુક્તિઓ ધ્યાનમાં લેવાં પડે
છે; $p$ વડે હકીકતમાં અસત્ય હોય એવું વાક્ય પ્રતીકિત કરીએ તોપણ
$p \Entails/ q$ છે.

છતાં લૂઈસના નીચેના વાક્ય જેવાં ``જો \dots તો\dots'' વાક્યોમાં
કંઈક વિચિત્ર લાગે છે:
\begin{quote}
જો ચંદ્ર લીલા પનીરનો બનેલો હોય, તો $2+2=4$.
\end{quote}
તેમજ નીચેના અનુમાનોમાં પણ:
\begin{quote}
  ચંદ્ર લીલા પનીરનો બનેલો નથી. તેથી, જો ચંદ્ર લીલા પનીરનો બનેલો
  હોય, તો $2+2=4$.

  $2+2 = 4$. તેથી, જો ચંદ્ર લીલા
  પનીરનો બનેલો હોય, તો $2+2=4$.
\end{quote}
પરંતુ ``જો \dots તો \dots''નો અર્થ માત્ર $\lif$ જ હોય, તો આ વાક્ય
કોઈ મુશ્કેલી વિના સત્ય અને આ અનુમાનો કોઈ મુશ્કેલી વિના માન્ય ગણાય.

બીજી ચિંતા પુનરુક્તિ $(!A \lif !B) \lor (!B \lif
!A)$માંથી ઊભી થાય છે. તે એવું સૂચવે છે કે અખબારમાંથી કોઈ પણ બે
નિવેદક વાક્યો~$S$ અને $T$ યાદૃચ્છિક રીતે લઈએ, તો ``જો $S$ તો
$T$, અથવા જો $T$ તો~$S$'' વાક્ય સત્ય હોવું જોઈએ.
% END OLP-0519
\endgroup


\begingroup
\makeatletter\def\ol@sectioncs{section}\makeatother

% BEGIN OLP-0520
\olfileid{cnt}{int}{str}

\olsection{કડક શરતી વિધાન}
\phantomsection\label{guunit:OLP-0520}


લૂઈસે \emph{કડક શરતી વિધાન}~$\strictif$ રજૂ કર્યું અને દલીલ કરી કે
ભૌતિક શરતી વિધાનને બદલે આ વિધાન ફલિતતાને અનુરૂપ છે. સત્યલક્ષી મોડલ
તર્કશાસ્ત્રમાં $!A \strictif !B$ને $\Box(!A \lif
!B)$ તરીકે વ્યાખ્યાયિત કરી શકાય છે. એટલે અનુરૂપ ભૌતિક શરતી
વિધાન અનિવાર્ય હોય ત્યારે અને માત્ર ત્યારે કડક શરતી વિધાન (કોઈ
જગતમાં) સત્ય હોય છે.

ભૌતિક શરતી વિધાનના વિરોધાભાસોની સામે કડક શરતી વિધાન કેટલું સફળ થાય
છે? અસત્ય પૂર્વપક્ષ ધરાવતું કડક શરતી વિધાન હોય કે સત્ય પરિણામ
ધરાવતું હોય, તે સત્ય પણ હોઈ શકે અને અસત્ય પણ. વધુમાં,
$(!A \strictif !B) \lor (!B \strictif !A)$ માન્ય નથી. કડક શરતી
વિધાન $!A \strictif !B$, $\lnot !A \lor !B$ને સમતુલ્ય પણ નથી;
એટલે તે સત્યતાફલનલક્ષી નથી.

અહીં \Log{S5}ના પરિપ્રેક્ષ્યમાં નીચેના સંબંધો મળે છે:
\sourcecorrection{OLMD-107}{આ ફલિતતા-સંબંધો તથા નીચેની અનિવાર્યતા-દલીલો દરેક મોડલ તંત્રમાં માન્ય નથી; અહીં S5નો પરિપ્રેક્ષ્ય સ્પષ્ટ કર્યો છે.}
\begin{align}
  !A \strictif !B & \Entails \lnot !A \lor !B
  \text{ પરંતુ:}\\
  \lnot !A \lor !B & \Entails/ !A \strictif !B\\
  !B & \Entails/ !A \strictif !B\\
  \lnot !A & \Entails/ !A \strictif !B\\
  \lnot(!A \strictif !B) & \Entails/ !A \land \lnot !B
  \text{ પરંતુ:}\\
  !A \land \lnot !B & \Entails \lnot(!A \strictif !B)
  \intertext{તેમ છતાં કડક શરતી વિધાન પર મોડસ પોનેન્સનો નિયમ લાગુ પડે છે:}
  !A, !A \strictif !B & \Entails !B
  \intertext{કડક શરતી વિધાન પરિણામોને એકત્ર કરે છે:}
  !A \strictif !B,  !A \strictif !C & \Entails !A \strictif (!B \land !C)
  \intertext{કડક શરતી વિધાનમાં પૂર્વપક્ષને મજબૂત કરી શકાય છે:}
  !A \strictif !B & \Entails (!A \land !C) \strictif !B
  \intertext{કડક શરતી વિધાન સંક્રામક પણ છે:}
  !A \strictif !B, !B \strictif !C & \Entails !A \strictif !C
  \intertext{છેવટે કડક શરતી વિધાન તેના પ્રતિપક્ષી સ્વરૂપ સાથે સમતુલ્ય છે:}
  !A \strictif !B & \Entails \lnot !B \strictif \lnot !A\\
  \lnot !B \strictif \lnot !A & \Entails !A \strictif !B
\end{align}
\sourcecorrection{OLMD-108}{મૂળમાં પાંચમી ફલિતતા-પંક્તિએ ભૌતિક સંયોજક મૂક્યો છે; તે રૂપ માટે દર્શાવેલી અ-ફલિતતા ખોટી છે, તેથી અહીં કડક સંયોજક મૂક્યો છે.}

\begin{prob}
  કડક શરતી વિધાન માટે માન્ય ન હોય એવા ફલિતતા-સંબંધોના
  \Log{S5}-પ્રત્યુદાહરણો આપો, એટલે કે નીચેનાં માટે:
  \begin{enumerate}
  \item $\lnot p \Entails/ \Box(p \lif q)$
  \item $q \Entails/ \Box(p \lif q)$
  \item $\lnot\Box(p \lif q) \Entails/ p \land \lnot q$
  \item $\Entails/ \Box(p \lif q) \lor \Box(q \lif p)$
  \end{enumerate}
\end{prob}

\begin{prob}
  માન્ય ફલિતતા-સંબંધો કડક શરતી વિધાન માટે સાચા છે તે નીચેના
  \Log{S5}-પ્રમાણો આપીને બતાવો:
  \begin{enumerate}
  \item  $\Box(!A \lif !B) \Entails \lnot !A \lor !B$
  \item $!A \land \lnot !B \Entails \lnot\Box(!A \lif !B)$
  \item $!A, \Box(!A \lif !B) \Entails !B$
  \item $\Box(!A \lif !B), \Box(!A \lif !C) \Entails \Box(!A \lif (!B
    \land !C))$
  \item $\Box(!A \lif !B) \Entails \Box((!A \land !C) \lif !B)$
  \item $\Box(!A \lif !B), \Box(!B \lif !C) \Entails \Box(!A
    \lif !C)$
  \item $\Box(!A \lif !B) \Entails \Box(\lnot !B \lif \lnot !A)$
  \item $\Box(\lnot !B \lif \lnot !A) \Entails \Box(!A \lif !B)$
  \end{enumerate}
  \end{prob}

તેમ છતાં કડક શરતી વિધાનના પોતાના ``વિરોધાભાસો'' છે. અસત્ય પૂર્વપક્ષ
અથવા સત્ય પરિણામ ધરાવતું ભૌતિક શરતી વિધાન સત્ય હોય છે; તેવી જ રીતે
\emph{અનિવાર્યપણે} અસત્ય પૂર્વપક્ષ અથવા અનિવાર્યપણે સત્ય પરિણામ ધરાવતું
કડક શરતી વિધાન સત્ય હોય છે. વધુમાં, \Log{S5}માં દરેક સત્ય કડક શરતી
વિધાન અનિવાર્યપણે સત્ય છે અને દરેક અસત્ય કડક શરતી વિધાન અનિવાર્યપણે
અસત્ય છે. બીજા શબ્દોમાં,
\begin{align}
  \Box \lnot !A & \Entails !A \strictif !B\\
  \Box !B & \Entails !A \strictif !B\\
  !A \strictif !B& \Entails \Box(!A \strictif !B)\\
  \lnot(!A \strictif !B) & \Entails \Box\lnot(!A \strictif !B)
\end{align}
$\strictif$ને ``ફલિત કરે છે'' એમ સમજીએ તો આ કોઈ સમસ્યા નથી.
આખરે તાર્કિક ફલિતતા-સંબંધો ગાણિતિક તથ્યો છે; તેથી તે સંજોગાધીન
હોઈ શકતા નથી. પરંતુ $\strictif$ને ``જો \dots તો \dots''નો અર્થ
પકડતો તાર્કિક સંયોજક માનવો હોય, તો ખાસ કરીને છેલ્લા બે સંબંધો
સમસ્યા ઊભી કરે છે. ``જો \dots તો \dots'' એવાં કેટલાંક વાક્યો
સંજોગાધીન રીતે સત્ય અથવા સંજોગાધીન રીતે અસત્ય હોય છે જ---હકીકતમાં,
સામાન્ય રીતે તે ન તો અનિવાર્ય હોય છે, ન તો અશક્ય.

\begin{prob}
  નીચેના \Log{S5}-પ્રમાણો આપો:
  \begin{enumerate}
    \item  $\Box \lnot !A \Entails !A \strictif !B$
    \item $!A \strictif !B \Entails \Box(!A \strictif !B)$
    \item $\lnot(!A \strictif !B)  \Entails \Box\lnot(!A \strictif !B)$
  \end{enumerate}
  તેના માટે $\strictif$ની વ્યાખ્યા વાપરો.
\end{prob}
% END OLP-0520
\endgroup


\begingroup
\makeatletter\def\ol@sectioncs{section}\makeatother

% BEGIN OLP-0521
\olfileid{cnt}{int}{cnt}

\olsection{પ્રતિતથ્યાત્મક શરતી વિધાનો}
\phantomsection\label{guunit:OLP-0521}


અંગ્રેજીમાં ``જો \dots તો \dots'' પ્રકારની વાક્યરચનાનું એક અત્યંત
સામાન્ય અને મહત્વનું સ્વરૂપ \emph{to be} ક્રિયાપદના ભૂતકાલીન
સંભાવનાર્થી રૂપથી બને છે: ``જો \dots એવું થયું હોત, તો \dots એવું
થયું હોત.'' સામાન્ય રીતે આવા શરતી વિધાનનો પૂર્વપક્ષ અસત્ય, એટલે
હકીકતથી વિપરીત, હોય છે; તેથી તેમને \emph{પ્રતિતથ્યાત્મક શરતી
  વિધાનો} કહે છે. \emph{to be}ના સંભાવનાર્થી રૂપને કારણે તેમને
\emph{સંભાવનાર્થી શરતી વિધાનો} પણ કહે છે. તે \emph{નિવેદક} શરતી
વિધાનોથી જુદાં છે, જેમનું સ્વરૂપ ``જો \dots એવું છે, તો \dots એવું
છે'' હોય છે. પ્રતિતથ્યાત્મક અને નિવેદક શરતી વિધાનોની સત્ય થવાની
શરતો જુદી છે. એડમ્સનું જાણીતું ઉદાહરણ જુઓ:
\begin{quote}
  જો ઓસ્વાલ્ડે કેનેડીની હત્યા કરી ન હતી, તો બીજા કોઈએ કરી હતી.
  
  જો ઓસ્વાલ્ડે કેનેડીની હત્યા કરી ન હોત, તો બીજા કોઈએ કરી હોત.
\end{quote}
પહેલું વિધાન નિવેદક છે અને બીજું પ્રતિતથ્યાત્મક. પહેલું સ્પષ્ટપણે
સત્ય છે: આપણે જાણીએ છીએ કે પ્રમુખ જૉન એફ. કેનેડીની હત્યા
\emph{કોઈએ} કરી હતી; અને જો એ વ્યક્તિ લી હાર્વી ઓસ્વાલ્ડ ન હતી
(વૉરેન અહેવાલથી વિપરીત), તો કેનેડીની હત્યા બીજા કોઈએ કરી હતી. બીજું
વિધાન કંઈક જુદું કહે છે. તે દાવો કરે છે કે જો ઓસ્વાલ્ડે કેનેડીની હત્યા
ન કરી હોત---એટલે ડૅલસમાં ગોળીબાર ટળી ગયો હોત અથવા નિષ્ફળ ગયો
હોત---તો ત્યાર પછી ઇતિહાસ એવી રીતે આગળ વધ્યો હોત કે બીજી કોઈ
હત્યાનો પ્રયાસ સફળ થયો હોત. તે સત્ય હોય, તો કેનેડીનું મૃત્યુ
સુનિશ્ચિત કરવા શક્તિશાળી પરિબળોએ કાવતરું રચ્યું હોવું જોઈએ (જેમ
કેનેડીની હત્યા અંગે કાવતરું થયું હતું એવો સિદ્ધાંત માનનારાઓમાંના
અનેક લોકો માને છે).

\emph{નિવેદક} શરતી વિધાનને ભૌતિક શરતી વિધાન યોગ્ય રીતે રજૂ કરે છે
કે નહિ, તે હજુ ચર્ચાનો વિષય છે. ખાસ કરીને, અંગ્રેજીનાં નિવેદક શરતી
વિધાનોની સત્ય થવાની શરતો ભૌતિક શરતી વિધાન આપે છે એ વિચાર સાથે
સુસંગત રીતે તેના વિરોધાભાસોને ``સમજાવી'' શકાય કે નહિ, તે પ્રશ્ન છે.
તેનાથી વિપરીત, પ્રતિતથ્યાત્મક શરતી વિધાનોનું ભૌતિક શરતી વિધાન વડે
યોગ્ય પ્રતીકીકરણ થઈ શકતું નથી, એ નિર્વિવાદ છે. આ સ્પષ્ટ છે, કારણ કે
પ્રતિતથ્યાત્મક વિધાનના પૂર્વપક્ષો સામાન્ય રીતે અસત્ય હોવા છતાં,
અસત્ય પૂર્વપક્ષ ધરાવતાં બધાં પ્રતિતથ્યાત્મક વિધાનો સત્ય હોતાં નથી.
દાખલા તરીકે, તમે વૉરેન અહેવાલ માનતા હો અને કેનેડીની હત્યા માટે કોઈ
કાવતરું ન રચાયું હોય, તો એડમ્સનું પ્રતિતથ્યાત્મક
વિધાન આવું ઉદાહરણ છે.

કારણસંબંધી તર્કવિચારમાં પ્રતિતથ્યાત્મક શરતી વિધાનો મહત્વની ભૂમિકા
ભજવે છે: કારણસંબંધ વ્યક્ત કરવો એ તેમનો મુખ્ય ઉપયોગ છે. દાખલા તરીકે,
દીવાસળી ઘસવાથી તે સળગે છે; તેને ``જો આ દીવાસળી ઘસાત, તો તે સળગી
ઊઠત'' એમ કહી શકાય. ભૌતિક શરતી વિધાનો, અને સામાન્ય રીતે નિવેદક
શરતી વિધાનો, આ સંબંધ વ્યક્ત કરી શકતા નથી: ``દીવાસળી ઘસવામાં આવે છે
$\lif$ દીવાસળી સળગે છે'' એ વિધાન દીવાસળી ક્યારેય ઘસવામાં ન આવે
તો સત્ય છે, ભલે ઘસવામાં આવે ત્યારે શું થાય તે ગમે તે હોય. એથી પણ
વધુ વિચિત્ર વાત એ છે કે ``દીવાસળી ઘસવામાં આવે છે $\lif$ દીવાસળી
ફૂલોનો ગુચ્છો બની જાય છે'' એ વિધાન પણ દીવાસળી ક્યારેય ઘસવામાં ન
આવે તો સત્ય છે; પરંતુ ઘસવામાં આવે તો તે ચોક્કસ ફૂલોનો ગુચ્છો નહિ
બને.

પ્રતિતથ્યાત્મક શરતી વિધાનોનું યોગ્ય તર્કશાસ્ત્ર ચોક્કસ શું છે, તેની
હજુ ચર્ચા થાય છે. સ્ટાલનેકર અને લૂઈસે તેમનું પ્રભાવશાળી વિશ્લેષણ
આપ્યું. તેમના મતે ``જો~$S$ એવું થયું હોત, તો~$T$ એવું થયું હોત''
પ્રતિતથ્યાત્મક વિધાન ત્યારે અને માત્ર ત્યારે સત્ય છે જ્યારે વાસ્તવિક
જગતની સ્થિતિને સૌથી વધુ મળતી આવતી અને જ્યાં~$S$ સત્ય છે એવી
પ્રતિતથ્યાત્મક સ્થિતિમાં (``શક્ય જગતમાં'') $T$ સત્ય હોય. આને
``સત્તાલક્ષી'' વિશ્લેષણ કહે છે, કારણ કે તે શક્ય જગતોની
અસ્તિત્વમીમાંસાનો આધાર લે છે. બીજાં વિશ્લેષણો શરતી સંભાવનાઓ અથવા
માન્યતાઓના પુનરીક્ષણના સિદ્ધાંતો વાપરે છે. પ્રતિતથ્યાત્મક વિધાનોના
જુદાં જુદાં સૂચિત તર્કશાસ્ત્રોની સંખ્યા ઘણી છે. લૂઈસ--સ્ટાલનેકરનું
પણ એકમાત્ર તર્કશાસ્ત્ર નથી: સ્ટાલનેકર અને લૂઈસે સૌથી નજીકનાં શક્ય
જગતોના આધારે મળતાં આવતાં વર્ણનો આપ્યાં હોવા છતાં, તેમની અલગ
પૂર્વધારણાઓમાંથી જુદાં માન્ય અનુમાનો ફલિત થાય છે.
% END OLP-0521
\endgroup


\OLEndChapterHook
% END OLP-0517
\endgroup


\begingroup
\makeatletter\def\ol@sectioncs{section}\makeatother

% BEGIN OLP-0522
\olchapter{cnt}{min}{ન્યૂનતમ ફેરફારનો અર્થવિચાર}
\olchapterid{cnt}{min}
\phantomsection\label{guunit:OLP-0522}


\begingroup
\makeatletter\def\ol@sectioncs{section}\makeatother

% BEGIN OLP-0523
\olfileid{cnt}{min}{int}

\olsection{પરિચય}
\phantomsection\label{guunit:OLP-0523}


સ્ટાલનેકર અને લૂઈસે ``જો દીવાસળી ઘસાત, તો તે સળગી ઊઠત'' જેવા
પ્રતિતથ્યાત્મક શરતી વિધાનોનું વિશ્લેષણ સૂચવ્યું. તેમનાં વિશ્લેષણો
આવાં વાક્યોની સત્ય થવાની શરતો યોગ્ય રીતે કેવી સમજવી તેની દરખાસ્તો
હતી. બંનેની પાછળનો વિચાર આ છે: પ્રતિતથ્યાત્મક શરતી વિધાન સત્ય છે
કે નહિ તે તપાસવા એવા શક્ય જગતો ધ્યાનમાં લેવાં પડે જે પૂર્વપક્ષને
સત્ય કરવા માટે વાસ્તવિક જગતથી ન્યૂનતમ પ્રમાણમાં જુદાં હોય. આ શક્ય
જગતોમાં પરિણામ સત્ય હોય, તો પ્રતિતથ્યાત્મક વિધાન સત્ય છે. દાખલા
તરીકે, ધારો કે મારા હાથમાં એક દીવાસળી અને દીવાસળીઓનું કાગળનું પૅકેટ
છે. વાસ્તવિક જગતમાં હું માત્ર તેમને જોઉં છું અને દીવાસળી ઘસું તો
શું થાય એ વિચારું છું. હું દીવાસળી ઘસું એવી, વાસ્તવિક જગતથી
ન્યૂનતમ ફેરફારવાળી સ્થિતિ એ છે જેમાં હું પગલું ભરવાનું નક્કી કરું
અને દીવાસળી ઘસું. આ ફેરફાર ન્યૂનતમ છે, કારણ કે બીજું કંઈ બદલાતું
નથી: હું સાથે હવામાં કૂદતો નથી; દીવાસળી ઘસતાં મારા વાળમાં પણ આગ
લાગતી નથી; મારી આંગળીઓની બધી તાકાત અચાનક જતી રહેતી નથી; અને
સુપર સોકર પાણીની પિસ્તોલના ઓચિંતા હુમલાથી હું એ જ સમયે ભીંજાઈ
જતો નથી, વગેરે. એ વૈકલ્પિક શક્યતામાં દીવાસળી સળગે છે. તેથી
``જો હું દીવાસળી ઘસત, તો તે સળગી ઊઠત'' એ સત્ય છે.

આ સહજ વિશ્લેષણને પ્રતિતથ્યાત્મક વિધાનોના તર્કશાસ્ત્ર માટેના ઔપચારિક
અર્થવિચાર સાથે જોડી શકાય છે. લૂઈસે પ્રતિતથ્યાત્મક સંયોજક માટે
``$\boxright$'' સંકેત રજૂ કર્યો અને સ્ટાલનેકરે~``$>$'' સંકેત વાપર્યો.
આપણે~$\cif$ વાપરીશું અને તેને વિધાનાત્મક તર્કશાસ્ત્રમાં દ્વિસ્થાની
સંયોજક તરીકે ઉમેરીશું. એટલે $!A \lif
!B$ સ્વરૂપનાં સૂત્રો ઉપરાંત~$!A \cif !B$ સ્વરૂપનાં
સૂત્રો પણ મળે છે. મોડલ તર્કશાસ્ત્રના સંબંધાત્મક અર્થવિચારની
જેમ, આ ઔપચારિક અર્થવિચાર એવા નિદર્શો પર આધારિત છે જેમાં
સૂત્રોનું મૂલ્યાંકન જગતો પર થાય છે. $\mSat{M}{!A \cif !B}[w]$
વ્યાખ્યાયિત કરતી સંતોષ-શરત કેટલાક (અન્ય) જગતો~$w'$ માટેના
$\mSat{M}{!A}[w']$ અને $\mSat{M}{!B}[w']$ વડે અપાય છે. કયા~$w'$?
સહજ રીતે કહીએ તો, $\mSat{M}{!A}[w']$ હોય એવાં~$w$ને સૌથી નજીકનાં
જગતો. આ માટે નિદર્શમાં ``નિકટતા''નો સંબંધ પણ સામેલ કરવો પડે છે.

લૂઈસે પ્રતિતથ્યાત્મક પરિસ્થિતિઓને આલેખમાં રજૂ કરવાની એક સમજણભરી
રીત રજૂ કરી. દરેક શક્ય જગત તેના આસપાસનાં અન્ય જગતોને સમાવતા
એકબીજાની અંદર આવેલા ગોળાઓના સમૂહના કેન્દ્રમાં હોય છે---આ ગોળાઓને
આપણે સમકેન્દ્રી વર્તુળો તરીકે દોરીએ છીએ. બે ગોળાઓ વચ્ચેના જગતો
કેન્દ્રના જગતથી એકસરખા નજીક હોય છે; અંદરના ગોળામાં આવેલાં જગતો
વધુ નજીક અને બહારના ગોળામાં આવેલાં જગતો વધુ દૂર હોય છે.
\begin{center}
\begin{tikzpicture}[scale=.7]\small
  \spheresystem{5}
  \draw (0,0) node {$w$};
  \propositionintersect{0}{4}{40}{3.5}
  {\tikzset{proposition/.append={smooth,tension=1.4}}
  \proposition[shift={(1.6,-1)}]{90}{1}{30}{4}}
  \path (1.5,3) node[below] {$\formula{B}$};
  \path (3,0) node {$\formula{A}$};
\end{tikzpicture}
\end{center}
સૌથી નજીકનાં $!A$-જગતો એવાં જગતો~$w'$ છે જ્યાં~$!A$ સંતોષાય છે
અને જે કેન્દ્રના જગત~$w$ની આસપાસના સૌથી નાના ગોળામાં (ભૂખરા
વિસ્તારમાં) આવેલાં છે. સહજ રીતે, સૌથી નજીકનાં બધાં $!A$-જગતોમાં
$!B$ સત્ય હોય, તો~$w$ પર $!A \cif !B$ સંતોષાય છે.
% END OLP-0523
\endgroup


\begingroup
\makeatletter\def\ol@sectioncs{section}\makeatother

% BEGIN OLP-0524
\olfileid{con}{min}{sph}

\olsection{ગોળા-નિદર્શો}
\phantomsection\label{guunit:OLP-0524}


પ્રતિતથ્યાત્મક વિધાનો માટે ઔપચારિક અર્થવિચાર આપવાનો એક રસ્તો લૂઈસના
અનૌપચારિક વિશ્લેષણને ગાણિતિક સંરચનામાં ફેરવવાનો છે. ત્યારે
જગત~$w$ની આસપાસના ગોળા જગતોના ગણો બને છે. ગોળા એકબીજાની અંદર
હોવાથી~$w$ની આસપાસના જગતોના આ ગણો ઉપગણ-સંબંધ મુજબ રેખીય ક્રમમાં
હોવા જોઈએ.

\begin{defn}
  \emph{ગોળા-નિદર્શ} એ ત્રિક~$\mModel{M} = \tuple{W, O, V}$ છે,
  જેમાં $W$ જગતોનો અખાલી ગણ છે, $V\colon \PVar \to \Pow{W}$
  મૂલ્યાંકન છે, અને $O\colon W \to \Pow{\Pow{W}}$ દરેક જગત~$w$ને
  \emph{ગોળાઓનું તંત્ર}~$O_w$ આપે છે. દરેક $w$ માટે $O_w$ જગતોના
  ગણોનો ગણ છે અને તેણે નીચેની શરતો સંતોષવી જોઈએ:
  \begin{enumerate}
  \item $O_w$નું \emph{કેન્દ્ર}~$w$ છે: $\{w\} \in O_w$.
  \item $O_w$ \emph{અંતર્વિષ્ટ} છે: જ્યારે પણ $S_1$, $S_2 \in O_w$
    હોય, ત્યારે $S_1 \subseteq S_2$ અથવા $S_2 \subseteq S_1$;
    એટલે $O_w$નો ક્રમ~$\subseteq$ મુજબ રેખીય છે.
  \item $O_w$ તેના ગોળાઓના અખાલી સંગ્રહોના યોગગણ હેઠળ બંધ છે.
  \item $O_w$ તેના ગોળાઓના અખાલી સંગ્રહોના છેદગણ હેઠળ બંધ છે.
  \end{enumerate}
\end{defn}

$O_w$ પાછળનો સહજ વિચાર એવો છે કે~$w$ની ``આસપાસનાં'' જગતો
~$w$થી કેટલાં દૂર છે તે મુજબ સ્તરોમાં ગોઠવાય છે. સૌથી અંદરનો
ગોળો એકલું~$w$ છે, એટલે ગણ~$\{w\}$: બીજા કોઈ પણ ગોળામાંનાં
જગતો કરતાં~$w$ પોતે~$w$ની વધુ નજીક છે. જો $S \subsetneq S'$, તો
$S' \setminus S$માંનાં જગતો~$S$માંનાં જગતો કરતાં~$w$થી વધુ દૂર
છે: $S' \setminus S$ એ~$S$ અને~$S'$ની બહારનાં જગતો વચ્ચેનો
``સ્તર'' છે. ખાસ કરીને, ગોળાની બાહ્ય સપાટીની અંદરનાં બધાં જગતો
તેમાં આવે છે એવું સમજવું જોઈએ; ગોળા માત્ર અલગ અલગ સ્તરો નથી.

\begin{figure}
\begin{center}
\begin{tikzpicture}[layerwidth=1.5,scale=.6]\tiny
  \spheresystem{3}
  \propositionintersect{0}{3}{60}{6}
  \path (0,0) node[world] {$w$};
  \spherepos{-30}{2}{node[world] {$w_2$}}
  \spherepos{220}{2}{node[world] {$w_3$}}
  \spherepos{90}{2}{node[world] {$w_1$}}
  \spherepos[shift={(.1,0)}]{10}{3}{node[world] {$w_5$}}
  \spherepos[shift={(.1,0)}]{-10}{3}{node[world] {$w_6$}}
  \spherepos{225}{3}{node[world] {$w_4$}}
  \spherepos{20}{4}{node[world] {$w_7$}}
  \path (5,0) node[left] {$p$};
\end{tikzpicture}
\caption{ગોળા-નિદર્શનો આલેખ}
\ollabel{fig:sphere-model}
\end{center}
\end{figure}

\olref{fig:sphere-model}નો આલેખ એવા ગોળા-નિદર્શને અનુરૂપ છે જેમાં
$W = \{w, w_1, \dots, w_7\}$ અને $V(p) = \{w_5, w_6,
w_7\}$. સૌથી અંદરનો ગોળો $S_1 = \{w\}$ છે. $w$થી સૌથી નજીકનાં
જગતો $w_1, w_2, w_3$ છે, તેથી તેનાથી મોટો આગળનો ગોળો
$S_2 = \{w, w_1, w_2, w_3\}$ છે. વધુ દૂરનાં જગતો $w_4$, $w_5$,
$w_6$ છે, તેથી સૌથી બહારનો ગોળો
$S_3 = \{w, w_1, \dots, w_6\}$ છે. $w$ની આસપાસના ગોળાઓનું તંત્ર
$O_w = \{S_1, S_2, S_3\}$ છે. જગત~$w_7$,~$w$ની આસપાસના કોઈ
પણ ગોળામાં નથી. જ્યાં $p$ સત્ય છે એવાં સૌથી નજીકનાં જગતો $w_5$
અને~$w_6$ છે; તેથી $p$ને સ્વીકારતો સૌથી નાનો ગોળો~$S_3$ છે.

ગોળા-નિદર્શ~$\mModel M$માં જગત~$w$ પર સૂત્ર~$!A$ના સંતોષને,
$\mSat{M}{!A}[w]$, વ્યાખ્યાયિત કરવા મોડલ સૂત્રો માટેની
વ્યાખ્યામાં $!B \cif !C$ની શરત ઉમેરીએ છીએ:

\begin{defn}
  $\mSat{M}{!B \cif !C}[w]$ ત્યારે અને માત્ર ત્યારે જ્યારે નીચેની
  બેમાંથી કોઈ એક શરત સંતોષાય:
  \begin{enumerate}
  \item\ollabel{sphere-vac} દરેક $u \in \bigcup O_w$ માટે $\mSat/{M}{!B}[u]$, અથવા
  \item\ollabel{sphere-nonvac} કોઈ $S \in O_w$ માટે,
    \begin{enumerate}
      \item કોઈ $u \in
        S$ માટે $\mSat{M}{!B}[u]$, અને
      \item દરેક $v \in S$ માટે કાં તો
        $\mSat/{M}{!B}[v]$ અથવા $\mSat{M}{!C}[v]$.
    \end{enumerate}
  \end{enumerate}
\end{defn}

આ વ્યાખ્યા મુજબ $\mSat{M}{!B \cif !C}[w]$ ત્યારે અને માત્ર ત્યારે
જ્યારે પૂર્વપક્ષ~$!B$,~$w$ની આસપાસના બધા ગોળાઓમાં બધે અસત્ય હોય,
અથવા કોઈ ગોળા~$S$માં $!B$ સત્ય હોય અને એ ``$!B$-સ્વીકારતા''
ગોળાનાં બધાં જગતોમાં ભૌતિક શરતી વિધાન $!B \lif !C$ સત્ય હોય.
ધ્યાન આપો કે સહજ સમજૂતી પરથી અપેક્ષા રાખી શકાય તેમ વ્યાખ્યામાં
$S$ સૌથી અંદરનો $!B$-સ્વીકારતો ગોળો હોવો જોઈએ એવી શરત મૂકી
નથી. પરંતુ~\olref{sphere-nonvac}ની શરત કોઈ ગોળા~$S$ માટે સાચી
હોય, તો~$S$માં સમાયેલો અને પૂર્વપક્ષને સ્વીકારતો દરેક ગોળો પણ એ
શરત સંતોષે છે; તેથી સૌથી અંદરનો એવો ગોળો અસ્તિત્વમાં
હોય તો તે પણ સંતોષે છે.
\sourcecorrection{OLMD-109}{મૂળમાં Sમાં સમાયેલા દરેક ગોળા માટે શરત સાચી કહેવાઈ છે; નાનામાં B-સાક્ષી ન પણ હોય. અહીં ફક્ત B-સ્વીકારતા નાના ગોળા માટેનો સાચો દાવો રાખ્યો છે.}

એ પણ ધ્યાન આપો કે ગોળા-નિદર્શની વ્યાખ્યા મુજબ સૌથી અંદરનો
$!B$-સ્વીકારતો ગોળો \emph{હોવો જ જોઈએ} એવું નથી: $!B$ને સ્વીકારતા
ગોળાઓની અનંત શ્રેણી $S_1 \supsetneq S_2 \supsetneq \dots \supsetneq
\{w\}$ હોઈ શકે, અને તેથી સૌથી અંદરનો $!B$-સ્વીકારતો ગોળો ન હોય.
એવા કિસ્સામાં કોઈ~$i$ માટે $S_i$, $S_{i+1}$, \dots{} એ બધાં ગોળાઓમાં
$!B \lif !C$ સત્ય હોય ત્યારે અને માત્ર ત્યારે
$\mSat{M}{!B \cif !C}[w]$ છે.
% END OLP-0524
\endgroup


\begingroup
\makeatletter\def\ol@sectioncs{section}\makeatother

% BEGIN OLP-0525
\olfileid{con}{min}{tf}

\olsection{પ્રતિતથ્યાત્મક વિધાનોની સત્યતા અને અસત્યતા}
\phantomsection\label{guunit:OLP-0525}


સૌથી નજીકનાં બધાં $!A$-જગતો $!B$-જગતો હોય, તો
$!A \cif !B$ પ્રતિતથ્યાત્મક વિધાન રિક્તતા વગર સત્ય હોય છે;
તે~\olref{fig:true}માં દર્શાવ્યું છે.
\begin{figure}
\begin{center}
\begin{tikzpicture}[scale=.7]\small
  \spheresystem{5}
  \draw (0,0) node {$w$};
  \propositionintersect{0}{4}{40}{3.5}
  {\tikzset{proposition/.append={smooth,tension=1.4}}
  \proposition[shift={(1.6,-1)}]{90}{1}{30}{4}}
  \path (1.5,3) node[below] {$\formula{B}$};
  \path (3,0) node {$\formula{A}$};
\end{tikzpicture}
\caption{રિક્તતા વગર સત્ય પ્રતિતથ્યાત્મક વિધાન}
\ollabel{fig:true}
\end{center}
\end{figure}
$w$ની આસપાસના ગોળાઓના તંત્રમાં $!A$ને સ્વીકારતો એક પણ ગોળો ન
હોય તોપણ પ્રતિતથ્યાત્મક વિધાન~$w$ પર સત્ય છે. એવા કિસ્સામાં તે
\emph{રિક્તતાને કારણે} સત્ય છે (જુઓ~\olref{fig:vacuous}).
\begin{figure}
\begin{center}
\begin{tikzpicture}[scale=.7]\small
  \spheresystem{5}
  \draw (0,0) node {$w$};
  \proposition{0}{6.5}{40}{3.5}
  {\tikzset{proposition/.append={smooth,tension=1.4}}
  \proposition[shift={(1.2,-1)}]{90}{1}{30}{4}}
  \path (1.5,3) node[below] {$\formula{B}$};
  \spherepos{0}{7}{node {$\formula{A}$}}
\end{tikzpicture}
\caption{રિક્તતાને કારણે સત્ય પ્રતિતથ્યાત્મક વિધાન}
\ollabel{fig:vacuous}
\end{center}
\end{figure}

તે બે રીતે અસત્ય હોઈ શકે છે. પહેલી રીતે, સૌથી નજીકનાં $!A$-જગતોમાં
કેટલાક $!B$-જગતો હોય પરંતુ બધાં ન હોય. આ કિસ્સામાં
$!A \cif \lnot !B$ પણ અસત્ય છે (જુઓ~\olref{fig:false}).
\begin{figure}
\begin{center}
\begin{tikzpicture}[scale=.7]\small
  \spheresystem{5}
  \draw (0,0) node {$w$};
  \propositionintersect{0}{4}{40}{3.5}
  {\tikzset{proposition/.append={smooth,tension=1.4}}
  \proposition[shift={(1.6,0)}]{90}{1}{30}{3}}
  \path (1.5,3) node[below] {$\formula{B}$};
  \path (3,0) node {$\formula{A}$};
\end{tikzpicture}
\caption{અસત્ય પ્રતિતથ્યાત્મક વિધાન; વિરુદ્ધ પરિણામવાળું પણ અસત્ય}
\ollabel{fig:false}
\end{center}
\end{figure}
સૌથી નજીકનાં $!A$-જગતો અને $!B$-જગતો એકબીજાને જરા પણ ન મળે,
તો $!A \cif !B$ અસત્ય છે. પરંતુ આ કિસ્સામાં સૌથી નજીકનાં બધાં
$!A$-જગતો $\lnot !B$-જગતો છે; તેથી $!A \cif \lnot !B$ સત્ય છે
(જુઓ~\olref{fig:false-opposite}).
\begin{figure}
\begin{center}
\begin{tikzpicture}[scale=.7]\small
  \spheresystem{5}
  \draw (0,0) node {$w$};
  \propositionintersect{0}{4}{40}{3.5}
  {\tikzset{proposition/.append={smooth,tension=1.4}}
  \proposition[shift={(-1,0)}]{90}{1}{30}{3}}
  \path (-1,3) node[below] {$\formula{B}$};
  \path (3,0) node {$\formula{A}$};
\end{tikzpicture}
\caption{અસત્ય પ્રતિતથ્યાત્મક વિધાન; વિરુદ્ધ પરિણામવાળું સત્ય}
\ollabel{fig:false-opposite}
\end{center}
\end{figure}

કડક શરતી વિધાનથી વિપરીત, પ્રતિતથ્યાત્મક વિધાનો સંજોગાધીન હોઈ શકે
છે. \olref{fig:contingent}માંનો ગોળા-નિદર્શ જુઓ. $u$થી સૌથી નજીકનાં
$!A$-જગતો બધાં $!B$-જગતો છે, તેથી $\mSat{M}{!A \cif
  !B}[u]$. પરંતુ~$v$થી સૌથી નજીકનાં કેટલાક $!A$-જગતો $!B$-જગતો
નથી, તેથી $\mSat/{M}{!A \cif !B}[v]$.
\begin{figure}
\begin{center}
\begin{tikzpicture}[scale=.6]\small
  \spheresystem{7}
  \spheresystem[shift={(0:3.1)},dashed]{4}
  \propositionintersect{45}{4}{40}{4.5}
  \begin{scope}
    \clip \propositionplot{45}{4}{40}{4.5} ;
    \spherelayer[shift={(0:3.1)}]{3}
  \end{scope}
  \proposition[shift={(1.4,-.5)}]{90}{1}{35}{5}
  \path (0,0) node {$u$};
  \path (0:3.1) node {$v$};
  \spherepos{35}{9}{node {$\formula{A}$}}
  \spherepos{80}{9}{node {$\formula{B}$}}
\end{tikzpicture}
\end{center}
\caption{સંજોગાધીન પ્રતિતથ્યાત્મક વિધાન}
\ollabel{fig:contingent}
\end{figure}
% END OLP-0525
\endgroup


\begingroup
\makeatletter\def\ol@sectioncs{section}\makeatother

% BEGIN OLP-0526
\olfileid{cnt}{min}{agg}

\olsection{પૂર્વપક્ષને મજબૂત કરવો}
\phantomsection\label{guunit:OLP-0526}


``પૂર્વપક્ષને મજબૂત કરવો'' એટલે
$!A \lif !C \Entails (!A \land !B) \lif !C$ એ અનુમાન.
ભૌતિક શરતી વિધાન માટે તે માન્ય છે, પરંતુ પ્રતિતથ્યાત્મક વિધાન માટે
અમાન્ય છે. ધારો કે આ દીવાસળી ઘસું તો તે સળગી ઊઠે, એ સાચું છે.
(એટલે દીવાસળી કે દીવાસળીઓના પૅકેટની ઘસવાની સપાટીમાં કોઈ ખામી નથી;
હું દીવાસળી તોડી નાખવાનો નથી, વગેરે.) પરંતુ આ દીવાસળી અવકાશમાં
ઘસું તો તે સળગી ઊઠે, એ સાચું નથી. તેથી નીચેનું અનુમાન અમાન્ય છે:
\begin{quote}
  જો દીવાસળી ઘસાત, તો તે સળગી ઊઠત.

  તેથી, જો દીવાસળી અવકાશમાં ઘસાત, તો તે સળગી ઊઠત.
\end{quote}

લૂઈસ--સ્ટાલનેકરનું શરતી વિધાનોનું વિશ્લેષણ આ સમજાવે છે: હું
અવકાશમાં દીવાસળી ઘસું એવું સૌથી નજીકનું જગત, હું દીવાસળી ઘસું
એવા સૌથી નજીકના જગત કરતાં વાસ્તવિક જગતથી ઘણું વધુ દૂર છે. તેથી
પછીના જગતમાં દીવાસળી સળગે છે, પરંતુ પહેલાંના જગતમાં નહિ. એવું જ
થવું જોઈએ.

\begin{ex}
  ગોળા-અર્થવિચાર આ અનુમાનને અમાન્ય ઠરાવે છે; એટલે $p
  \cif r \Entails/ (p \land q) \cif r$. નિદર્શ $\mModel{M}
  = \tuple{W, O, V}$ લો જેમાં $W = \{w, w_1, w_2\}$,
  $O_w = \{\{w\}, \{w, w_1\}, \{w, w_1, w_2\}\}$,
  $V(p) = \{w_1, w_2\}$, $V(q) = \{w_2\}$ અને
  $V(r) = \{w_1\}$. $p$ને સ્વીકારતો ગોળો $S = \{w,
  w_1\}$ છે અને તેના બધાં જગતોમાં $p \lif r$ સત્ય છે; તેથી
  $\mSat{M}{p \cif r}[w]$. વળી $(p \land q)$ને સ્વીકારતો ગોળો
  $S' = \{w, w_1, w_2\}$ છે, પરંતુ
  $\mSat/{M}{(p \land q) \lif r}[w_2]$; તેથી
  $\mSat/{M}{(p \land q) \cif r}[w]$ (જુઓ~\olref{fig:antecedent-strengthening}).

\begin{figure}
\begin{center}
\begin{tikzpicture}[layerwidth=1.5,scale=.6]\tiny
  \spheresystem{3}
  \propositionintersect{-10}{3}{30}{5}
  \begin{scope}
    \clip plot[smooth,tension=1] coordinates
          {(0,4.5) (30:.95) (4.5,-3.5)} -- (4.5,4.5) ;
    \spherefill{2}
  \end{scope}
  \draw[proposition] plot[dashed,smooth,tension=1] coordinates
          {(0,4.5) (30:.95) (4.5,-3.5)} ;
  \proposition{30}{2}{30}{5}
  \path (0,0) node[world] {$w$};
  \spherepos{30}{2}{node[world] {$w_1$}}
  \spherepos{-10}{3}{node[world] {$w_2$}}
  \spherepos{-10}{4}{node {$q$}}
  \spherepos{30}{4}{node {$r$}}
  \draw (1,4.5) node[below] {$p$};
\end{tikzpicture}
\caption{પૂર્વપક્ષને મજબૂત કરવાના નિયમનું પ્રત્યુદાહરણ}
\ollabel{fig:antecedent-strengthening}
\end{center}
\end{figure}
\end{ex}
% END OLP-0526
\endgroup


\begingroup
\makeatletter\def\ol@sectioncs{section}\makeatother

% BEGIN OLP-0527
\olfileid{cnt}{min}{tra}

\olsection{સંક્રામકતા}
\phantomsection\label{guunit:OLP-0527}


ભૌતિક શરતી વિધાન માટે શૃંખલા-નિયમ માન્ય છે: $!A \lif !B, !B
\lif !C \Entails !A \lif !C$. બીજા શબ્દોમાં, ભૌતિક શરતી વિધાન
સંક્રામક છે. શું પ્રતિતથ્યાત્મક વિધાનો માટે પણ એવું જ છે? સ્ટાલનેકરે
આપેલું નીચેનું ઉદાહરણ જુઓ.
\begin{quote}
  જો જે.~એડગર હૂવર રશિયન તરીકે જન્મ્યા હોત, તો તેઓ કોમ્યુનિસ્ટ બન્યા હોત.

  જો જે.~એડગર હૂવર કોમ્યુનિસ્ટ હોત, તો તેઓ દેશદ્રોહી હોત.

  તેથી, જો જે.~એડગર હૂવર રશિયન તરીકે જન્મ્યા હોત, તો તેઓ દેશદ્રોહી હોત.
\end{quote}
હૂવર વાસ્તવમાં જે સમયે જન્મ્યા એ જ સમયે અમેરિકાને બદલે રશિયામાં
જન્મ્યા હોત, તો તેઓ સોવિયેત સંઘમાં ઊછરીને કોમ્યુનિસ્ટ બન્યા હોત
(એવું ધારી લઈએ). તેથી પહેલો પૂર્વાધાર સત્ય છે. એ જ રીતે, બીજા
પૂર્વાધારને અલગથી લઈએ તો તે પણ સત્ય છે. પરંતુ તારણ અસત્ય છે:
હૂવર સોવિયેત સંઘમાં જન્મ્યા હોત, તો સંભવતઃ તેઓ કટ્ટર કોમ્યુનિસ્ટ
બન્યા હોત અને (પોતાના દેશના) દેશદ્રોહી ન બન્યા હોત. સત્યમૂલ્યોનું
આ સહજ નિયોજન સ્ટાલનેકર--લૂઈસના વિશ્લેષણથી પણ સમર્થિત થાય છે.
હૂવરનું જન્મસ્થળ બદલાય એટલો જ ફેરફાર ધરાવતું આપણા જગતથી સૌથી
નજીકનું શક્ય જગત એ છે જેમાં હૂવર સોવિયેત સંઘના સારા નાગરિક તરીકે
ઊછરે છે. પહેલો પૂર્વાધાર અને તારણ, બંનેનો પૂર્વપક્ષ એ જગતમાં સત્ય
છે; અને ત્યાં હૂવર કોમ્યુનિસ્ટ પક્ષના વફાદાર સભ્ય છે, તેથી દેશદ્રોહી
નથી. પરંતુ બીજો પૂર્વાધાર મૂલવવા જુદું જગત જોવું પડે: હૂવર
કોમ્યુનિસ્ટ હોય એવું સૌથી નજીકનું જગત. એમાં તેઓ અમેરિકામાં જન્મ્યા,
પછી કોમ્યુનિસ્ટ બન્યા અને તેથી દેશદ્રોહી બન્યા.\footnote{અલબત્ત, આ
  ઉદાહરણનું બળ સમજવા કેટલીક અધિભૌતિક અને રાજકીય ધારણાઓ સ્વીકારવી
  પડે; દાખલા તરીકે, હૂવર રશિયન માતાપિતાને જન્મ્યા હોત એ શક્યતા,
  અથવા 1950ના દાયકામાં અમેરિકાના કોમ્યુનિસ્ટો પોતાના દેશના
  દેશદ્રોહી હતા એ ધારણા.}

\begin{prob}
  પ્રતિતથ્યાત્મક વિધાનોની સંક્રામકતા નિષ્ફળ જાય તેનું વિશ્વસનીય,
  સહજ ઉદાહરણ શોધો.
\end{prob}

\begin{ex}\ollabel{ex:trans-counterex}
  ગોળા-અર્થવિચાર આ અનુમાનને અમાન્ય ઠરાવે છે; એટલે $p
  \cif q, q \cif r \Entails/ p \cif r$. નિદર્શ $\mModel{M}
  = \tuple{W, O, V}$ લો જેમાં $W = \{w, w_1, w_2\}$,
  $O_w = \{\{w\}, \{w, w_1\}, \{w, w_1, w_2\}\}$,
  $V(p) = \{w_2\}$, $V(q) = \{w_1, w_2\}$ અને
  $V(r) = \{w_1\}$. $p$ને સ્વીકારતો ગોળો $S = \{w, w_1,
  w_2\}$ છે અને તેમાંનાં બધાં જગતોમાં $p \lif q$ સત્ય છે; તેથી
  $\mSat{M}{p \cif q}[w]$. વળી $q$ને સ્વીકારતો ગોળો
  $S' = \{w, w_1\}$ છે અને તેમાંનાં બધાં જગતોમાં
  $q \lif r$ સત્ય છે; તેથી
  $\mSat{M}{q \cif r}[w]$. પરંતુ $p$ને સ્વીકારતા ગોળા
  $\{w, w_1, w_2\}$માં જગત~$w_2$ છે, જ્યાં
  $\mSat/{M}{p \lif r}[w_2]$.
  \sourcecorrection{OLMD-110}{મૂળમાં $q\lif r$ બધા જગતોમાં સત્ય કહેતાં ખોટું અસંતોષ-સૂત્ર મૂકાયું છે; અહીં તે જ ભૌતિક શરતી વિધાન બધા જગતોમાં સત્ય હોવાનું સ્પષ્ટ લખ્યું છે.}
\end{ex}

\begin{prob}
\olref[cnt][min][tra]{ex:trans-counterex}ના પ્રત્યુદાહરણને અનુરૂપ
ગોળાઓનો આલેખ દોરો.
\end{prob}

\begin{prob}
  \olref[cnt][min][tra]{ex:trans-counterex}માં જગત $w_2$ પર હૂવર
  રશિયામાં જન્મ્યા છે, કોમ્યુનિસ્ટ છે અને દેશદ્રોહી નથી; જગત $w_1$
  પર તેઓ અમેરિકામાં જન્મ્યા છે, કોમ્યુનિસ્ટ છે અને દેશદ્રોહી છે.
  આ નિદર્શમાં $w_1$,~$w$ની $w_2$ કરતાં વધુ નજીક છે. શું આ જરૂરી
  છે? રશિયામાં જન્મવું, હૂવર કોમ્યુનિસ્ટ બને તેની તુલનામાં વધુ દૂરની
  શક્યતા છે એવી ધારણા વગરનું પ્રત્યુદાહરણ આપી શકો છો?
\end{prob}
% END OLP-0527
\endgroup


\begingroup
\makeatletter\def\ol@sectioncs{section}\makeatother

% BEGIN OLP-0528
\olfileid{cnt}{min}{cpo}

\olsection{પ્રતિપક્ષી રૂપાંતર}
\phantomsection\label{guunit:OLP-0528}


ભૌતિક અને કડક શરતી વિધાનો પોતાનાં પ્રતિપક્ષી સ્વરૂપો સાથે સમતુલ્ય
છે. પ્રતિતથ્યાત્મક વિધાનો એવાં નથી. ક્રાત્ઝરે આપેલું ઉદાહરણ જુઓ:
\begin{quote}
  જો ગોથે 1832માં મૃત્યુ પામ્યા ન હોત, તો તેઓ અત્યારે પણ મૃત્યુ પામેલા હોત.

  જો ગોથે અત્યારે મૃત્યુ પામેલા ન હોત, તો તેઓ 1832માં મૃત્યુ પામ્યા હોત.
\end{quote}
પહેલું વાક્ય સત્ય છે: મનુષ્યો સેંકડો વર્ષ જીવતા નથી. બીજું વાક્ય
સ્પષ્ટપણે અસત્ય છે: ગોથે અત્યારે મૃત્યુ પામેલા ન હોત, તો હજુ જીવતા
હોત અને તેથી 1832માં મૃત્યુ પામ્યા ન હોઈ શકે.

\begin{ex}\ollabel{ex:contraposition-counterex}
  ગોળા-અર્થવિચાર પ્રતિપક્ષી રૂપાંતરને અમાન્ય ઠરાવે છે; એટલે $p
  \cif q \Entails/ \lnot q \cif \lnot p$. $p$ને ``ગોથે 1832માં
  મૃત્યુ પામ્યા ન હતા'' અને $q$ને ``ગોથે અત્યારે મૃત્યુ પામેલા છે''
  તરીકે સમજો. તેને નિદર્શ $\mModel{M_1} = \tuple{W, O, V}$માં
  રજૂ કરી શકાય છે, જ્યાં $W = \{w, w_1, w_2\}$,
  $O_w = \{\{w\}, \{w, w_1\}, \{w, w_1, w_2\}\}$,
  $V(p) = \{w_1, w_2\}$ અને $V(q) = \{w, w_1\}$. એટલે $w$
  વાસ્તવિક જગત છે જેમાં ગોથે 1832માં મૃત્યુ પામ્યા અને હજુ મૃત્યુ
  પામેલા છે; $w_1$ નજીકનું જગત છે જેમાં તેઓ, માનો કે, 1833માં
  મૃત્યુ પામ્યા અને હજુ મૃત્યુ પામેલા છે; $w_2$ દૂરનું જગત છે જેમાં
  તેઓ હજુ જીવતા છે.
\begin{figure}
\begin{center}
\begin{tikzpicture}[scale=.6,layerwidth=1.5]\tiny
  \spheresystem{3}
  \begin{scope}
    \clip \propositionplot{0}{.8}{50}{4.5} -- (4.5,-4.5) -- (-4.5,-4.5) -- (-4.5,4.5) -- (4.5,4.5);
    \spherelayer{3}
  \end{scope}
  \propositionintersect{0}{2}{110}{5.5}
  \proposition{0}{.8}{50}{4.5}
  \path (0,0) node[world] {$w$};
  \spherepos{0}{2}{node[world] {$w_1$}}
  \spherepos{-50}{3}{node[world] {$w_2$}}
  \spherepos{28}{3}{node {$q$}}
  \spherepos{42}{3}{node {$\lnot q$}}
  \spherepos{-55}{4}{node {$p$}}
  \spherepos{-67}{4}{node {$\lnot p$}}
\end{tikzpicture}
\caption{પ્રતિપક્ષી રૂપાંતરનું પ્રત્યુદાહરણ}
\ollabel{fig:contraposition}
\end{center}
\end{figure}
$p$ને સ્વીકારતો ગોળો $S = \{w, w_1\}$ છે અને તેમાંનાં બધાં
જગતોમાં $p \lif q$ સત્ય છે; તેથી $\mSat{M_1}{p \cif q}[w]$.
પરંતુ $\lnot q$ને સ્વીકારતા ગોળા $\{w, w_1, w_2\}$માં જગત~$w_2$
છે, જ્યાં $q$ અસત્ય અને $p$ સત્ય છે; તેથી
$\mSat/{M_1}{\lnot q \lif \lnot p}[w_2]$.
\sourcecorrection{OLMD-111}{મૂળમાં જગતનું ગોળા-તંત્ર $O$ને બદલે $O_w$ હોવું જોઈએ અને નિદર્શનું નામ $M_1$ આપ્યા પછી સંતોષ-સૂત્રોમાં M છે; અહીં બંને સંજ્ઞાઓ સુસંગત કરી છે.}
\end{ex}
% END OLP-0528
\endgroup


\OLEndChapterHook
% END OLP-0522
\endgroup


\OLEndPartHook
% END OLP-0516
